%% ma: L11-B04-0026
%% lop: 11
%% chuong: 1
%% bai: 4
%% dang: 11.4.12
%% loai: DS
%% muc: [NB, NB, NB, VD]
%% dap_an: "ĐĐĐS"
%% nguon: "de-on-tap-giua-ki-1-toan-11-nam-2026-2027-truong-thpt-ngo-quyen-da-nang.pdf (M09, Đề ôn tập giữa kì 1 Toán 11 THPT Ngô Quyền, Đà Nẵng), Mã đề 0114, Phần II Câu 1"
%% kien_thuc: "Bài 3 (hàm số lượng giác), Bài 4 (phương trình lượng giác cơ bản)"
%% trang_thai: chinh-thuc
%% ghi_chu: "Đã duyệt 29/10/2026."
\begin{ex}
Cho hàm số $f(x)=\tan2x-1$. Khi đó:
\choiceTF{\True Giá trị của hàm số tại $x=\dfrac{\pi}{8}$ bằng $0$}{\True Giá trị của hàm số tại $x=\dfrac{\pi}{3}$ bằng $-\sqrt3-1$}{\True Hàm số đã cho là một hàm số tuần hoàn}{Có đúng $3$ giá trị $x$ thuộc khoảng $(0;\pi)$ thỏa mãn hàm số đạt giá trị bằng $-2$}
\loigiai{a) Đúng: $\tan\dfrac\pi4-1=0$.

b) Đúng: $\tan\dfrac{2\pi}3-1=-\sqrt3-1$.

c) Đúng: hàm số tuần hoàn với chu kì $\dfrac\pi2$.

d) Sai: $\tan2x=-1\Leftrightarrow x=-\dfrac\pi8+\dfrac{k\pi}2$; trên $(0;\pi)$ có $2$ giá trị $\dfrac{3\pi}8$, $\dfrac{7\pi}8$.}
\end{ex}
